% 第3章 研究方法
% 対応する草稿: research/thesis-drafts/chapter3-methods-draft.md

\subsection{実験の全体設計}\label{subsec:experiment-design}

% TODO: 3つの実験（基本calibration, プロンプト比較, 日英比較）の概要

\subsection{データセット}\label{subsec:dataset}

% TODO: 4ドメイン × 各25問 = 100問
% 出典、構築方法、品質管理

\subsection{対象モデル}\label{subsec:models}

% TODO: GPT-5, Claude Sonnet 4.6, Gemini 2.5, Swallow
% バージョン情報、実行条件

\subsection{プロンプト設計}\label{subsec:prompts}

% TODO: 4方式（0-100整数, 5段階, %, 自然言語）
% プロンプトテンプレートを verbatim で掲載

\subsection{評価指標}\label{subsec:metrics}

% TODO: ECE, Brier Score, AUROC の数式定義
% DEL 公理との対応表

\subsection{実験手順}\label{subsec:procedure}

% TODO: API呼び出し、正誤判定、3回繰り返し平均
